\chapter{결론}\label{chap:conclusion}

본 논문은 SSD의 처리량 포화를 고려하여 읽을 행을 선택하는 문제를 다시 정의하고, 그 최적해를 구하는 Global DP와 계산이 가벼운 선택 방법인 Tiles를 제시하였다.

Tiles는 Neuron Chunking의 greedy 방법과 비슷한 읽기 시간을 훨씬 짧은 선택 시간으로 얻었다. $896\times896$ 가중치 행렬의 50\% 예산에서 두 방법의 SSD 읽기 시간은 0.4 ms 안팎으로 비슷하였으나, 선택 시간은 Tiles가 0.01 ms, greedy 방법이 1.04 ms였다. 그 결과 선택과 읽기를 합한 전체 시간은 Tiles가 약 3분의 1이었고, 네 가지 shape 모두에서 Tiles가 더 짧았다. 그 대신 Tiles의 보존 중요도는 이 조건에서 3.6\%p 낮았다.

비용 모델에서의 비교에서도 두 방법은 비슷하였다. 같은 조건에서 Tiles는 Global DP가 구한 최적 비율의 95.9\%에, greedy 방법은 96.8\%에 도달하였다.

본 논문의 한계와 향후 과제는 다음과 같다.

\begin{itemize}
\item \textbf{모델의 정확도}: 선택 방법을 보존 중요도로 평가하였다. 선택한 행만 읽어 추론한 모델의 정확도는 다루지 않았다.
\item \textbf{모델과 장치}: 중요도는 언어 모델 하나에서 수집하였고, 모든 측정은 노트북 한 대에서 수행하였다. Neuron Chunking이 대상으로 한 VLM과 다른 저장 장치에서의 확인이 필요하다.
\item \textbf{비용 모델}: Two-line 모델은 포화 이후의 처리량을 일정하다고 둔다. 실측 처리량에는 요청 크기에 따른 국소적인 변동이 남아 있었고, 이 제약을 두지 않은 hinge 모델의 test 오차가 더 작았다.
\end{itemize}
